\frontchapter{Abstract}

{\singlespacing\noindent
\textbf{PERALTA, DAVE CHRISTIAN P.} Institute of Computer Science, College of
Arts and Sciences, University of the Philippines Los Baños. September 2026.
\textbf{PathFinder: An Indoor Navigation Web Application for the Physical
Sciences Building in the University of the Philippines Los Baños.}\par}

{\singlespacing\noindent
Adviser: Fermin Roberto G. Lapitan\par}

Indoor wayfinding remains challenging in multistory buildings, where
satellite-based positioning is unreliable and indoor positioning systems require
costly, hardware-dependent infrastructure that is difficult to scale across
floors and buildings. This study developed and evaluated a web-based indoor
navigation application for the Physical Sciences Building at the University of
the Philippines Los Baños, designed as a scalable and software-driven
alternative to hardware-based positioning and static physical floor plans. The
application provides a navigation interface that guides users between selected
origin and destination points, and a layout builder that allows building
administrators to create and publish digital map layouts for display within the
application. Usability was evaluated with two participant groups: six regular
users, who assessed navigation accuracy through a standardized usability
questionnaire and open-ended questions, and four building administrators, who
assessed the layout builder's functionality qualitatively. Users reported a mean
usability score of 81.67 out of 100, rated ``Excellent,'' with positive feedback
on route clarity, while administrators found the layout builder intuitive but
identified specific areas for refinement. These findings support the application
as a viable, scalable approach to indoor navigation.

\noindent
\textbf{Keywords:} indoor navigation, wayfinding, layout builder, web
application
